%! The Victory of Orthogonality — Unit 7, Lesson 7.4 — Slides (Beamer)
\documentclass[aspectratio=169, 11pt]{beamer}
\usepackage{linalg-beamer}   % provides \burgundyheader{} and \sectionlabel[color]{}
\usetikzlibrary{arrows.meta, positioning, calc}
\newcommand{\uu}{\mathbf{u}}
\newcommand{\vv}{\mathbf{v}}
\newcommand{\xx}{\mathbf{x}}
\newcommand{\qq}{\mathbf{q}}
\newcommand{\zero}{\mathbf{0}}
\newcommand{\T}{^{\top}}

\begin{document}

% TITLE SLIDE (CourseName is not defined in beamer, so the name is literal)
{
\setbeamercolor{background canvas}{bg=burgundy}
\begin{frame}[plain]
  \vspace{2.8cm}\hspace{0.55cm}%
  \begin{minipage}{0.9\textwidth}
    {\color{goldacc}\bfseries\small LESSON 7.4}\\[0.5cm]
    {\color{white}\bfseries\LARGE The Victory of Orthogonality}\\[1.0cm]
    {\color{blushmid}\small Linear Algebra~$\cdot$~Unit 7: The Singular Value Decomposition}
  \end{minipage}
\end{frame}
}

% HOOK
\begin{frame}[t, plain]
  \burgundyheader{What idea has been doing the work all unit?}
  \sectionlabel{Hook}
  \begin{columns}[T]
    \begin{column}{0.52\textwidth}
      Singular vectors came in \textbf{perpendicular} pairs. Compression and PCA both used perpendicular
      directions. Symmetric matrices had perpendicular eigenvectors. Least squares was a perpendicular
      projection.\\[6pt]
      One idea runs through all of it: \textbf{orthogonality}.
    \end{column}
    \begin{column}{0.44\textwidth}
      \begin{block}{The question}
        What is so special about a matrix whose columns are \textbf{perpendicular unit vectors}?
      \end{block}
    \end{column}
  \end{columns}
  \vspace{2pt}
  \begin{block}{Today}
    Meet the \textbf{orthogonal matrices} --- and see why the SVD writes \emph{every} matrix using them.
  \end{block}
\end{frame}

% ORTHOGONAL MATRIX + FREE INVERSE
\begin{frame}[t, plain]
  \burgundyheader{Orthogonal matrices: one test, one free inverse}
  \sectionlabel{Skill}
  \begin{itemize}
    \setlength{\itemsep}{6pt}
    \item $Q$ is \textbf{orthogonal} when its columns are perpendicular unit vectors --- one equation:
          \[ Q\T Q=I \qquad Q=\tfrac1{\sqrt2}\begin{bsmallmatrix} 1 & 1\\ 1 & -1 \end{bsmallmatrix}
             \ \Rightarrow\ Q\T Q=\begin{bsmallmatrix} 1 & 0\\ 0 & 1 \end{bsmallmatrix}. \]
    \item Diagonal $1$'s $=$ unit length; off-diagonal $0$'s $=$ perpendicular.
    \item \textbf{The inverse is free:} $Q^{-1}=Q\T$ --- no elimination, no determinant.
  \end{itemize}
\end{frame}

% LENGTH PRESERVED
\begin{frame}[t, plain]
  \burgundyheader{Orthogonal matrices preserve length}
  \sectionlabel[goldacc]{Idea}
  \begin{columns}[T]
    \begin{column}{0.54\textwidth}
      \[ |Q\xx|^2=(Q\xx)\T(Q\xx)=\xx\T Q\T Q\,\xx=\xx\T\xx=|\xx|^2. \]
      So $|Q\xx|=|\xx|$ for \textbf{every} $\xx$: an orthogonal matrix moves vectors \textbf{rigidly} --- a
      rotation or a reflection, no stretching.
    \end{column}
    \begin{column}{0.42\textwidth}
      \begin{block}{Check: $\xx=(3,1)$}
        $Q\xx=\tfrac1{\sqrt2}(4,2)$, and
        \[ |Q\xx|^2=\tfrac12(16+4)=10=|\xx|^2. \]
        Same length.
      \end{block}
    \end{column}
  \end{columns}
\end{frame}

% ROTATE-STRETCH-ROTATE
\begin{frame}[t, plain]
  \burgundyheader{The SVD is rotate--stretch--rotate}
  \sectionlabel{Skill}
  \begin{columns}[T]
    \begin{column}{0.52\textwidth}
      $A=U\Sigma V\T$ with $U,V$ orthogonal, $\Sigma$ diagonal. Right to left:
      \begin{itemize}
        \setlength{\itemsep}{4pt}
        \item $V\T$ --- \textbf{turns} the input to the axes;
        \item $\Sigma$ --- \textbf{stretches} by $\sigma_1,\sigma_2$;
        \item $U$ --- \textbf{turns} to the output.
      \end{itemize}
      \vspace{4pt}
      All the stretching lives in $\Sigma$; the orthogonal factors only rotate.
    \end{column}
    \begin{column}{0.46\textwidth}
      \centering
      \begin{tikzpicture}[font=\scriptsize, node distance=5.5mm,
          box/.style={draw=burgundy, fill=blush, rounded corners=1pt, minimum width=8mm, minimum height=6mm, align=center},
          ar/.style={-{Latex[length=1.6mm]}, burgundy, thick}]
        \node[box] (x)  {$\xx$};
        \node[box, right=of x] (v)  {$V\T$};
        \node[box, right=of v] (s)  {$\Sigma$};
        \node[box, right=of s] (u)  {$U$};
        \node[box, right=of u] (ax) {$A\xx$};
        \draw[ar] (x)--(v);
        \draw[ar] (v)--(s);
        \draw[ar] (s)--(u);
        \draw[ar] (u)--(ax);
        \node[burgundy] at ($(v)+(0,0.8)$) {turn};
        \node[burgundy] at ($(s)+(0,0.8)$) {stretch};
        \node[burgundy] at ($(u)+(0,0.8)$) {turn};
      \end{tikzpicture}
    \end{column}
  \end{columns}
\end{frame}

% THE VICTORY
\begin{frame}[t, plain]
  \burgundyheader{The victory --- and the revolution}
  \sectionlabel[goldacc]{Payoff}
  \begin{itemize}
    \setlength{\itemsep}{6pt}
    \item Orthogonal $U,V$ are \textbf{reversible for free}, \textbf{length-preserving}, and \textbf{never
          amplify errors} --- so modern algorithms compute with them (the ``revolution'').
    \item Orthogonality is the thread through the course: projections \& least squares (U4), perpendicular
          eigenvectors (U6), perpendicular singular vectors for \emph{every} matrix (U7).
  \end{itemize}
  \vspace{2pt}
  \begin{block}{Unit 7 in one line}
    Every matrix is \textbf{rotate--stretch--rotate}: two perfect orthogonal turns and one diagonal stretch.
    That is orthogonality's victory --- the reason the SVD, compression, and PCA all work.
  \end{block}
\end{frame}

\end{document}
